\documentclass{axiom}
\usepackage{axiom-macros}
\title{Functions}
\author{}
\date{April 2026}
\begin{document}
\maketitle
\source{Spivak, \textit{Calculus}, 4th Edition, Chapter 3}

\begin{dependencies}
  \dep{foundations/01-field-axioms}{Field axioms --- arithmetic on function values}
  \dep{foundations/02-order-axioms}{Order axioms --- inequality on \(\R\)}
\end{dependencies}

\section{The Definition of a Function}
\noindent We now introduce four definitions that formalize the notion of a function.

\noindent Vi introduserer nå fire definisjoner som formaliserer funksjonsbegrepet.

\begin{CJK}{UTF8}{min}
\noindent 我々は今、関数の概念を形式化する四つの定義を導入する。
\end{CJK}

\begin{definition}[Function]\label{d:function}
\end{definition}

\begin{definition}[Domain]\label{d:domain}
\end{definition}

\begin{definition}[Function Value]\label{d:function-value}
\end{definition}

\begin{definition}[Equality of Functions]\label{d:function-equality}
\end{definition}

\newpage
\section{Derived Definitions}
\noindent We now define several specific functions and derived operations on functions.

\noindent Vi definerer nå flere spesifikke funksjoner og avledede operasjoner på funksjoner.

\begin{CJK}{UTF8}{min}
\noindent ここで、複数の具体的な関数と関数に対する派生演算を定義する。
\end{CJK}

\subsection{Pointwise Arithmetic}

\begin{definition}[Sum]\label{d:sum}
\end{definition}

\begin{definition}[Difference]\label{d:difference}
\end{definition}

\begin{definition}[Product]\label{d:product}
\end{definition}

\begin{definition}[Quotient]\label{d:quotient}
\end{definition}

\subsection{Composition}

\begin{definition}[Composition]\label{d:composition}
\end{definition}

\subsection{Specific Functions}

\begin{definition}[Identity Function]\label{d:identity}
\end{definition}

\begin{definition}[Characteristic Function]\label{d:characteristic}
\end{definition}

\newpage
\section{Consequences}
\noindent We now derive three consequences of composition.

\noindent Vi utleder nå tre konsekvenser av komposisjon.

\begin{CJK}{UTF8}{min}
\noindent 以下では合成の三つの帰結を導く。
\end{CJK}

\begin{proposition}[Associativity of Composition]\label{p:compose-assoc}
\end{proposition}

\begin{proof}
\end{proof}

\begin{proposition}[Identity is Neutral for Composition]\label{p:compose-identity}
\end{proposition}

\begin{proof}
\end{proof}

\begin{proposition}[Non-Commutativity of Composition]\label{p:compose-noncomm}
\end{proposition}

\begin{proof}
\end{proof}

\end{document}
